\section{INTRODUCTION}

Contact-rich manipulation requires understanding both the global scene and the
local contact interface. An external RGB camera provides object identity,
orientation context, and coarse spatial layout, but is occluded precisely where
the gripper meets the object. A vision-based tactile sensor captures
high-resolution contact geometry---depth and surface normals at sub-millimeter
resolution---yet its field of view spans only the gel surface (${\sim}12$\,mm
after cropping). Neither modality alone is sufficient: tactile images of
geometrically symmetric objects (e.g.\ a circle, a square) are rotationally
ambiguous, while RGB lacks the contact detail needed for dense surface
reconstruction.

Prior work on visuo-tactile fusion typically assumes pixel-aligned image pairs
or operates separate, differently-sized models per modality configuration,
making it difficult to isolate the benefit of fusion from capacity
differences. Moreover, most methods share a single encoder for both
modalities, overlooking the growing availability of modality-specific
pretrained representations---tactile foundation models such as
T3~\cite{zhao2024transferable} and Sparsh~\cite{higuera2024sparsh}, and
self-supervised visual encoders such as MAE~\cite{he2022mae} and
DINOv2~\cite{oquab2023dinov2}---that offer stronger inductive biases when
matched to their respective domains.

\method{} addresses these gaps with three design decisions.
%
First, each modality is encoded by a \emph{frozen, modality-specific} encoder
selected through systematic ablation (T3 for tactile, MAE ViT-L for RGB);
freezing anchors sim-to-real transfer and guarantees that no encoder capacity
shifts between modality configurations.
%
Second, an \emph{asymmetric bottleneck fusion trunk} treats tactile tokens as
spatial queries and RGB as read-only key/value context, funnelled through a
small set of learnable bottleneck tokens. Because the output dimensionality
follows the query count, the decoder input is invariant to which modalities
are present.
%
Third, the dense prediction (DPT) and pose estimation paths are
\emph{architecturally decoupled}: the DPT decoder taps the tactile encoder
directly at its native 1024-dim multiscale features and trains on every step,
while modality dropout is confined to the pose path through the 768-dim
fusion trunk. This prevents dropout-induced dilution of dense supervision.

A single model trained with modality dropout supports three inference
modes---RGB+tactile, tactile-only, and RGB-only---with identical trainable
parameters, enabling a fair ablation of each modality's contribution.
Co-trained on simulated and real tactile data at an empirically optimized
ratio, \method{} transfers to real sensors without fine-tuning.

The contributions of this work are:
\begin{itemize}
\item An asymmetric bottleneck-attention architecture with decoupled dense and
      pose paths for fusing non-aligned RGB and tactile images, producing
      per-pixel depth, surface normals, and SE(2) pose in the tactile frame;
\item A systematic encoder ablation across six pretrained backbones, showing
      that modality-matched frozen encoders (T3 + MAE) outperform a single
      shared encoder;
\item A fair single-model ablation protocol---identical trainable parameters
      and decoder input across three input configurations---that isolates each
      modality's contribution and reveals complementary roles (tactile
      $\rightarrow$ dense geometry, RGB $\rightarrow$ pose disambiguation);
\item Sim-to-real transfer via frozen features and sim+real co-training,
      validated on real vision-based tactile sensors. \todo{sensor name}
\end{itemize}
